% 两层随机性分别作用于训练输出和固定输出上的新观测。
\begin{tikzpicture}[figure picture,foundation picture]
 \path[use as bounding box] (0,7) rectangle (112,72);
 \node[foundation input,fill=FigureBlueFill,minimum width=30mm] (p) at (54,65) {未知分布$\mathcal P$};
 \node[foundation input,fill=FigureBlueFill,minimum width=28mm] (s) at (20,42) {训练集$S\sim\mathcal P^n$};
 \node[foundation input,fill=FigureBlueFill,minimum width=28mm] (z) at (88,42) {新观测$z\sim\mathcal P$};
 \node[foundation model,minimum width=28mm] (a) at (20,20) {学习算法$A$};
 \node[foundation output,fill=FigureYellowFill,minimum width=30mm] (h) at (54,20) {$\hat h=A(S)$};
 \node[foundation loss,fill=FigureRedFill,minimum width=30mm] (l) at (88,20) {$\ell(\hat h,z)$};
 \draw[foundation flow] (p)--(s); \draw[foundation flow] (p)--(z);
 \draw[foundation flow] (s)--(a); \draw[foundation flow] (a)--(h);
 \draw[foundation flow] (z)--(l);
 \draw[foundation flow] (h.east)--(l.west);
 \draw[foundation frame,dashed] (3,11) rectangle (71,51);
\end{tikzpicture}%
